% GENERATED FILE. Edit canonical lessons, then run npm run generate:books.
% canonical-source-hash: fnv1a64:520dbb95adf38d64
% canonical-lessons: ML-C51-gratitude, ML-C51-favour, ML-C51-respect, ML-C51-blessing, ML-C51-salutation

\chapter{Courtesy}
\label{ch:ml-courtesy}
\hlchaptermodality{51}

\begin{chapteropening}
\textbf{By the end of this chapter:} \emph{I can thank, respect and bless in the formal Malayalam register, and say how each word differs from its everyday partner.}

Complete ML-C51-salutation and say all five words in order.
\end{chapteropening}

\section[kṛtajñata]{\ml{കൃതജ്ഞത} (kṛtajñata) — gratitude}

\label{lesson:ML-C51-gratitude}



\begin{quote}
Before the new run of words: what did \emph{puṟappeṭuka} mean, and what did \emph{viṭa} mean?
\end{quote}

\subsection*{You'll want to know: \ml{കൃതജ്ഞത}}
\textbf{\ml{കൃതജ്ഞത}} (\emph{kṛtajñata}) — \textquotedblleft{}gratitude\textquotedblright{}.

Sanskrit again, and legible once opened: \emph{kṛta} is 'done', \emph{jña} is 'knowing'. The grateful person is the one who knows what was done for them.

Malayalam already has \ml{നന്ദി} (\emph{nandi}) for the everyday thank-you, and that one is the family's own. So the pair here is the pattern this book keeps meeting: a native word for the ordinary moment and a Sanskrit one for the formal register, the letter, the speech.

Kannada has \ka{ಕೃತಜ್ಞತೆ} (\emph{kṛtajñate}), the same borrowing with a Kannada ending. Where the formal register is concerned the sisters agree far more often than they do in the kitchen.

\begin{grammarlens}[title={What you've built}]
The first of five courtesies.
\end{grammarlens}

\subsection*{Your turn}
Say these aloud:

\begin{itemize}
\raggedright
  \item \emph{kṛtajñata}
  \item it once more, slowly
  \item \emph{kṛtajñata}, then \emph{nandi}, the formal and the everyday
\end{itemize}

\begin{tcolorbox}[breakable,colback=teal!4,colframe=teal!35!black,title={Before you move on}]
What does \ml{കൃതജ്ഞത} mean? (\textquotedblleft{}gratitude\textquotedblright{}.) Say it once more.
\end{tcolorbox}
\section[upakāraṁ]{\ml{ഉപകാരം} (upakāraṁ) — a favour, a good turn}

\label{lesson:ML-C51-favour}



\begin{quote}
Before the new one: what did \emph{kṛtajñata} mean?
\end{quote}

\subsection*{You'll want to know: \ml{ഉപകാരം}}
\textbf{\ml{ഉപകാരം}} (\emph{upakāraṁ}) — \textquotedblleft{}a favour, a good turn\textquotedblright{}.

\emph{Upa} is 'near, alongside'; \emph{kāra} is 'a doing'. A favour is a doing-alongside — help brought up close.

The word carries a polite adverb with it in daily speech, an 'as a kindness' shaped much like the \ml{ദയവായി} this book taught for 'please'. Both are courtesies built by turning a noun sideways into an adverb, which is a habit worth noticing rather than learning one word at a time.

Kannada has \ka{ಉಪಕಾರ} (\emph{upakāra}), the same word once more.

\begin{grammarlens}[title={What you've built}]
Two: gratitude, and a favour.
\end{grammarlens}

\subsection*{Your turn}
Say these aloud:

\begin{itemize}
\raggedright
  \item \emph{upakāraṁ}
  \item it once more, slowly
  \item \emph{upakāraṁ}, then \emph{dayavāyi}, two courtesies made the same way
\end{itemize}

\begin{tcolorbox}[breakable,colback=teal!4,colframe=teal!35!black,title={Before you move on}]
What does \ml{ഉപകാരം} mean? (\textquotedblleft{}a favour, a good turn\textquotedblright{}.) Say it once more.
\end{tcolorbox}
\section[bahumānaṁ]{\ml{ബഹുമാനം} (bahumānaṁ) — respect}

\label{lesson:ML-C51-respect}



\begin{quote}
Before the new one: what did \emph{upakāraṁ} mean?
\end{quote}

\subsection*{You'll want to know: \ml{ബഹുമാനം}}
\textbf{\ml{ബഹുമാനം}} (\emph{bahumānaṁ}) — \textquotedblleft{}respect\textquotedblright{}.

\emph{Bahu} is 'much'; \emph{māna} is 'regard, esteem'. To respect is to regard much, and the Sanskrit says so plainly.

That \emph{māna} runs back to a root about thinking, and the same root gives Sanskrit its word for the mind. Malayalam took that mind-word too: it is sitting inside \ml{മനസ്സിലാക്കുക}, the 'put it in the mind' this book used for understanding. Respect and understanding share an ancestor.

Kannada does not use this word for respect; it says \ka{ಗೌರವ} (\emph{gaurava}), which is built on 'heavy'. Regarding much, or weighing much — the two sisters chose different metaphors off the same shelf.

\begin{grammarlens}[title={What you've built}]
Three. Gratitude, a favour, respect.
\end{grammarlens}

\subsection*{Your turn}
Say these aloud:

\begin{itemize}
\raggedright
  \item \emph{bahumānaṁ}
  \item it once more, slowly
  \item \emph{bahumānaṁ}, then \emph{manassilākkuka}, and find the piece they share
\end{itemize}

\begin{tcolorbox}[breakable,colback=teal!4,colframe=teal!35!black,title={Before you move on}]
What does \ml{ബഹുമാനം} mean? (\textquotedblleft{}respect\textquotedblright{}.) Say it once more.
\end{tcolorbox}
\section[anugrahaṁ]{\ml{അനുഗ്രഹം} (anugrahaṁ) — a blessing}

\label{lesson:ML-C51-blessing}



\begin{quote}
Before the new one: what did \emph{bahumānaṁ} mean?
\end{quote}

\subsection*{You'll want to know: \ml{അനുഗ്രഹം}}
\textbf{\ml{അനുഗ്രഹം}} (\emph{anugrahaṁ}) — \textquotedblleft{}a blessing\textquotedblright{}.

\emph{Anu} is 'after, along'; \emph{graha} is 'a taking hold'. A blessing is a taking-hold on someone's behalf — favour laid on with a hand.

That \emph{graha} is a busy word. \ml{ഗ്രഹം} in Malayalam is a planet, and it is the same noun: in the older astronomy a planet was a seizer, something that took hold of a life. One root, a blessing at one end and a planet at the other.

Kannada blesses with \ka{ಆಶೀರ್ವಾದ} (\emph{āśīrvāda}), a different Sanskrit compound built on speech rather than on grasp. A blessing said, against a blessing laid on.

\begin{grammarlens}[title={What you've built}]
Four. Gratitude, a favour, respect, a blessing.
\end{grammarlens}

\subsection*{Your turn}
Say these aloud:

\begin{itemize}
\raggedright
  \item \emph{anugrahaṁ}
  \item it once more, slowly
  \item \emph{anugrahaṁ}, then \emph{bahumānaṁ}, both of them things given rather than taken
\end{itemize}

\begin{tcolorbox}[breakable,colback=teal!4,colframe=teal!35!black,title={Before you move on}]
What does \ml{അനുഗ്രഹം} mean? (\textquotedblleft{}a blessing\textquotedblright{}.) Say it once more.
\end{tcolorbox}
\section[vandanaṁ]{\ml{വന്ദനം} (vandanaṁ) — a salutation}

\label{lesson:ML-C51-salutation}



\begin{quote}
Before the new one: what did \emph{anugrahaṁ} mean?
\end{quote}

\subsection*{You'll want to know: \ml{വന്ദനം}}
\textbf{\ml{വന്ദനം}} (\emph{vandanaṁ}) — \textquotedblleft{}a salutation\textquotedblright{}.

Off a Sanskrit root about praising and saluting. \ml{വന്ദനം} is the act of paying respects, and it is the word on a printed invitation rather than in a doorway.

The doorway word this book taught first was \ml{നമസ്കാരം}, and its own sense is a making of a bow. So Malayalam holds two Sanskrit nouns for what is nearly one physical act: the bend of the body, and the honour paid by it.

Kannada has \ka{ವಂದನೆ} (\emph{vandane}), and its own everyday greeting is built on that first word as well. The sisters furnished their doorways from the same shop.

\begin{grammarlens}[title={What you've built}]
Five, and the run is closed: gratitude, a favour, respect, a blessing, a salutation.
\end{grammarlens}

\subsection*{Your turn}
Say these aloud:

\begin{itemize}
\raggedright
  \item \emph{vandanaṁ}
  \item it once more, slowly
  \item all five in order — \emph{kṛtajñata}, \emph{upakāraṁ}, \emph{bahumānaṁ}, \emph{anugrahaṁ}, \emph{vandanaṁ}
\end{itemize}

\begin{tcolorbox}[breakable,colback=teal!4,colframe=teal!35!black,title={Before you move on}]
What does \ml{വന്ദനം} mean? (\textquotedblleft{}a salutation\textquotedblright{}.) Say it once more.
\end{tcolorbox}
